\documentclass[11pt,a4paper]{article}

% ---------- Layout ----------
\usepackage[margin=1in]{geometry}
\usepackage{microtype}

% ---------- Math ----------
\usepackage{amsmath,amssymb,amsthm,mathtools}

% ---------- Tables, links, bibliography ----------
\usepackage{booktabs}
\usepackage{graphicx}
\usepackage{hyperref}
\hypersetup{colorlinks=true,linkcolor=blue!60!black,urlcolor=blue!60!black,citecolor=blue!60!black}
\usepackage[round,authoryear]{natbib}

% ---------- Title ----------
\title{Fitting a Hyperspherical-Harmonic Curve to a Software-Skill Corpus:\\
  {\large A Spherical-Manifold Extension of the Learned-Latent-Curve Family}}

\author{Shant Tchatalbachian}

\date{August 5, 2026}

\begin{document}

\maketitle

\begin{abstract}
We present a curve-fitting method whose parameter manifold is the $N$-sphere $S^N$ rather than the $[0,1]^2$ torus used by the prior learned-latent-curve family. The model uses the canonical hyperspherical-harmonic basis $\{Y^{S^N}_{\ell,m}\}$ --- eigenfunctions of the Laplace--Beltrami operator on $S^N$ with eigenvalues $-\ell(\ell+N-1)$ --- composed with a learned M\"obius transformation $\phi_\theta \in \mathrm{PSL}(2,\mathbb{C})$ that reparameterizes the domain. The composition is novel at the application layer: we did not identify a prior work combining hyperspherical-harmonic basis, learned M\"obius reparameterization, and corpus-fit evaluation on a software-skill corpus. We evaluate the model on two splits of a software-skill corpus ($N=49$ and $N=70$) against a capacity-matched flat Fourier baseline. On both splits, the spherical model achieves a higher holdout $R^2$ than the flat baseline at fewer parameters (matched-parameter ablation). The absolute holdout $R^2$ is positive at $N=49$ ($+0.379$) and negative at $N=70$ ($-0.236$), indicating that the corpus is too small for the spherical model to outperform predicting the corpus mean at the larger split. We describe the method, the dataset, the evaluation protocol, and the limitations that bound the result.
\end{abstract}

\section{Introduction}

\subsection{Problem statement}

Given $N$ items in a $K$-dimensional feature space, find (i) a scalar ordering coordinate $t_i \in [0,1]$ for each item and (ii) a fixed-width $D$-dimensional embedding $z_i \in \mathbb{R}^D$ such that interpolation between adjacent items in the ordering is well defined and sparse regions of the ordering are detectable. This arises in corpus-audit tasks: a practitioner needs to know which parts of an ordering are thin and need more work.

\subsection{Background: the learned-latent-curve family}

The learned-latent-curve family re-expands a 1-D coordinate $t$ into a $D$-dimensional embedding via a Fourier basis with $k$ shared learned frequencies:
\begin{equation}
z_j(t) = a_{j,0} + \sum_{m=1}^{k} \big( a_{j,m} \sin(2\pi f_m t) + b_{j,m} \cos(2\pi f_m t) \big).
\label{eq:flat}
\end{equation}
With $k$ frequencies and $D$ output dimensions, the model has $P = k + D(1+2k)$ parameters; at $D=384$, $k=8$, $P = 6{,}536$. The closed-form coefficient solve at fixed frequencies is the standard Tikhonov ridge
\begin{equation}
C^\star = (\Phi^\top \Phi + \lambda I)^{-1} \Phi^\top Z,
\label{eq:ridge}
\end{equation}
the sanity floor for any gradient fit.

The flat $[0,1]^2$ parameter manifold used by the original formulation has zero scalar curvature, trivial topology, and trivial holonomy. Boundary behaviour is a free choice (periodic, zero-padded, or clamped). For corpora whose intrinsic geometry is well-modelled by a curved manifold, this is a poor inductive bias.

\subsection{This paper}

We extend the curve family to the $N$-sphere $S^N$ (default $N=2$, the Riemann sphere $\widehat{\mathbb{C}} \cong \mathbb{CP}^1$). The parameter manifold becomes $S^N$; the basis becomes the canonical hyperspherical-harmonic basis $Y^{S^N}_{\ell,m}$ (eigenfunctions of $\Delta_{S^N}$); and the domain is reparameterized by a learned M\"obius transformation $\phi_\theta \in \mathrm{PSL}(2,\mathbb{C})$. We evaluate on a 49-item and a 70-item split of a software-skill corpus against a capacity-matched flat Fourier baseline.

\section{Notation}

\begin{itemize}
\item $D$ --- embedding dimension (default 384).
\item $N$ --- corpus size; $N_{\text{train}}, N_{\text{holdout}}$ --- split sizes.
\item $k$ --- number of Fourier frequencies in the flat baseline (default 8).
\item $L$ --- maximum harmonic degree in the spherical basis (default 3).
\item $\mathbf{x} \in S^N \subset \mathbb{R}^{N+1}$ --- input point on the sphere.
\item $Y^{S^N}_{\ell,m}(\mathbf{x})$ --- real spherical harmonic of degree $\ell$, order $m$.
\item $\phi_\theta(z) = (az+b)/(cz+d)$ --- M\"obius transformation with $ad-bc=+1$, parameterised by $\theta$.
\item $\chi(z_1,z_2;z_3,z_4) = (z_1-z_3)(z_2-z_4)/((z_1-z_4)(z_2-z_3))$ --- cross-ratio.
\end{itemize}

\section{The Hyperspherical-Harmonic Curve}

\subsection{Model}

For $\mathbf{x} \in S^N \subset \mathbb{R}^{N+1}$,
\begin{equation}
\gamma(\mathbf{x}) = \mathbf{b} + \sum_{\ell=0}^{L} \sum_{m} \mathbf{a}_{\ell,m}\, Y^{S^N}_{\ell,m}\!\big(\phi_\theta(\mathbf{x})\big),
\label{eq:hyperspherical}
\end{equation}
where $\mathbf{b} \in \mathbb{R}^D$ is the per-dim bias, $\mathbf{a}_{\ell,m} \in \mathbb{R}^D$ are the per-dim per-basis coefficients, and $\phi_\theta$ is the M\"obius reparameterization.

The parameter count is $D \cdot n_{\text{basis}} + D + n_{\text{m\"obius}}$, where $n_{\text{basis}} = \sum_{\ell=0}^{L}(2\ell+1)$ (at $L=3$ on $S^2$, $n_{\text{basis}}=16$) and $n_{\text{m\"obius}} = 8$ raw real components minus 1 constraint ($ad-bc=+1$) minus 1 $\mathrm{PSL}(2,\mathbb{C})$ identification = 6 real degrees of freedom. At $D=384$, $L=3$, $N=2$: $P = 384 \cdot 16 + 384 + 6 = 6{,}534$.

\subsection{Real spherical harmonics}

The real spherical harmonics $Y^{S^2}_{\ell,m}$ on $S^2$ are constructed via explicit Legendre polynomials and the cos/sin split:
\begin{align}
Y^c_{\ell,0}(\theta,\phi) &= N_\ell^0 \cdot P_\ell^0(\cos\theta) \\
Y^c_{\ell,m}(\theta,\phi) &= \sqrt{2} \cdot N_\ell^m \cdot P_\ell^m(\cos\theta) \cdot \cos(m\phi) \quad (m > 0) \\
Y^c_{\ell,-m}(\theta,\phi) &= \sqrt{2} \cdot N_\ell^m \cdot P_\ell^m(\cos\theta) \cdot \sin(m\phi) \quad (m > 0)
\end{align}
with $N_\ell^m = \sqrt{(2\ell+1)/(4\pi) \cdot (\ell-m)!/(\ell+m)!}$ the standard normalisation. At $L=3$ on $S^2$ the basis has $1 + 3 + 5 + 7 = 16$ functions; at $L=3$ on $S^3$ the basis has $30$ functions.

\subsection{M\"obius reparameterization}

$\phi_\theta(z) = (az+b)/(cz+d)$ with $a,b,c,d \in \mathbb{C}$ and $ad-bc=+1$. Eight stored real components, six degrees of freedom after the $ad-bc=+1$ constraint and the $\mathrm{PSL}(2,\mathbb{C})$ identification: $\mathrm{Re}(a), \mathrm{Im}(a), \mathrm{Re}(b), \mathrm{Im}(b), \mathrm{Re}(c), \mathrm{Im}(c), \mathrm{Re}(d), \mathrm{Im}(d)$. We refine $\theta$ via L-BFGS-B with the closed-form ridge (Eq.~\ref{eq:ridge}) re-solved at each step. Because every $\varphi \in \mathrm{PSL}(2,\mathbb{C})$ preserves the cross-ratio identically by definition, the cross-ratio check verifies implementation consistency: it fails if $\varphi_\theta$ is miscoded, not if the learned map is a poor fit.

\section{Dataset and Evaluation Protocol}

\subsection{Dataset}

The corpus consists of software-skill specifications (engineering artifacts). Each item has a 9-dimensional binary feature vector recording which of nine primitive capabilities the skill implements. We consider two splits:
\begin{itemize}
\item \textbf{Split A} (49 items): the first 49 items alphabetically. Train $N_{\text{train}}=35$, holdout $N_{\text{holdout}}=14$.
\item \textbf{Split B} (70 items): all items. Train $N_{\text{train}}=49$, holdout $N_{\text{holdout}}=21$.
\end{itemize}

The 9-D binary feature space has principal-component concentration $\mathrm{PC1}+\mathrm{PC2} = 0.652$ at Split A and $0.548$ at Split B. The split sizes are chosen because (i) a curve is the honest model for small $N$ and (ii) the $S^2/L=3$ basis needs at least $N \geq 48$ items to fit.

\subsection{Baseline}

The capacity-matched baseline is a flat Fourier curve on $[0,1]^2$ with $k=2$ (the 2-D tensor-product extension of Eq.~\ref{eq:flat}), giving $5 \times 5 = 25$ basis functions and $9{,}984$ parameters. Both models are fitted with the closed-form ridge (Eq.~\ref{eq:ridge}) and the same ridge regularisation $\lambda$.

\subsection{Evaluation metric}

The headline metric is \emph{matched-parameter ablation}: holdout $R^2$ at fewer or equal parameters, with the same split and the same ridge regularisation. We report the absolute holdout $R^2$ (not just the delta) so the result is interpretable against the corpus-mean baseline ($R^2 = 0$).

\subsection{Reproducibility}

The basis construction is deterministic given $(\ell, m)$. The M\"obius refinement uses L-BFGS-B; the initial point is $\theta_0$ with $a=d=1$, $b=c=0$ (identity M\"obius). The ridge $\lambda$ is fixed across both models. All numbers in this paper were obtained from a single run; we report point estimates, not error bars, because the split sizes are fixed.

\section{Results}

\subsection{Split A (49 items)}

\textbf{Hyperspherical model}: $R^2_{\text{holdout}} = +0.379$ (6,534 parameters). \\
\textbf{Flat Fourier baseline}: $R^2_{\text{holdout}} = -0.359$ (9,984 parameters). \\
\textbf{Delta} (spherical $-$ flat): $+0.737$.

The hyperspherical model outperforms the flat baseline at fewer parameters. The absolute $R^2$ is positive on this split.

\subsection{Split B (70 items)}

\textbf{Hyperspherical model}: $R^2_{\text{holdout}} = -0.236$ (6,534 parameters). \\
\textbf{Flat Fourier baseline}: $R^2_{\text{holdout}} = -0.756$ (9,984 parameters). \\
\textbf{Delta} (spherical $-$ flat): $+0.520$.

Both models perform worse than predicting the corpus mean ($R^2 = 0$) on this split. The hyperspherical model is less negative than the flat baseline; the relative comparison is positive. The absolute result is negative.

\subsection{Calibration checks}

\begin{itemize}
\item Spectral mass $\rho = \sum_{\ell \geq 1} \|a_{:,\ell}\|^2 / \sum_{\ell \geq 0} \|a_{:,\ell}\|^2 \geq 0.10$: measured $0.977$ (A) / $0.983$ (B).
\item High-degree mass $\sum_{\ell > L/2} \|a_{:,\ell}\|^2 / \mathrm{total} \leq 0.40$: measured $0.206$ (A) / $0.178$ (B).
\item Cross-ratio check on 100 held-out 4-tuples: max residual $3.08 \times 10^{-7}$. This is consistent with float64 noise; every $\mathrm{PSL}(2,\mathbb{C})$ element preserves $\chi$ exactly, so the residual measures implementation precision, not fit quality.
\item M\"obius refinement (train $R^2$): identity $0.9125 \to$ refined $0.9211$, $\Delta = +0.009$. Train-only; no holdout effect measured.
\end{itemize}

\section{Discussion}

\subsection{What this result does and does not show}

The relative comparison (spherical vs flat on the same split) is positive on both splits. The absolute $R^2$ is positive only on Split A; on Split B, both models are worse than the corpus mean. The result supports the claim that, on this corpus, a spherical parameter manifold is a better inductive bias than a flat one; it does not support the claim that either model is a good fit in absolute terms.

\subsection{Why a sphere?}

The flat $[0,1]^2$ manifold has zero scalar curvature, trivial topology, and trivial holonomy. The Riemann sphere $S^2$ has constant positive scalar curvature $+2$, simply-connected topology ($\chi = 2$), and non-trivial holonomy. For corpora whose intrinsic geometry is well-modelled by $S^2$, the curved parameter manifold is a better inductive bias.

\subsection{Why M\"obius?}

$\mathrm{PSL}(2,\mathbb{C})$ is the full automorphism group of $\widehat{\mathbb{C}}$ and the 6-parameter family of invertible conformal domain warps on the Riemann sphere. The M\"obius refinement is not isometric: only the subgroup $\mathrm{PSU}(2) \cong SO(3)$ preserves the round metric. This non-isometry is what makes the refinement able to change the fit.

\subsection{Limitations}

\begin{enumerate}
\item The implementation uses explicit Legendre + cos/sin split. Production could swap to a maintained basis library.
\item We evaluate on a single target (binary primitive coverage). The result may not transfer to other targets (raw content, sentence-transformer embeddings).
\item Split B's negative absolute $R^2$ indicates the corpus is at the edge of the curve model's honest range. We do not know where that edge is.
\item $S^3/L=3$ (30 basis functions) would need $N \geq 90$ items. The current corpus is below that threshold.
\end{enumerate}

\section{Related Work}

We did not identify a prior work that combines hyperspherical-harmonic basis, learned M\"obius reparameterization, and corpus-fit evaluation on a software-skill corpus. We searched for related work using 11 keyword patterns across the closest candidates:

\textbf{arXiv:2601.20528 (Durastanti, 2026)}: Spectral Bayesian Regression on the Sphere. Covers Fourier-on-$S^2$ as Bayesian regression with statistical-theory focus; posterior contraction rates, Laplace--Beltrami smoothing splines. Does not cover: corpus fit, learned M\"obius reparameterization, sparse-cell detection.

\textbf{OpenReview g6UqpVislvH (ICLR 2022 submission)}: Generalized Fourier Features for Coordinate-Based Learning on Manifolds. Covers positional encoding for NeRF / panorama / SO(3) using spherical harmonics + SO(2)/SO(3) rotation shifts. Does not cover: corpus fit, $\mathrm{PSL}(2,\mathbb{C})$ M\"obius, sparse-cell detection.

The application-layer composition (hyperspherical basis + M\"obius + corpus fit) appears to be a gap in the published literature as of this writing.

\section{Conclusion}

The hyperspherical-harmonic curve is a candidate Stage-1 method for the learned-latent-curve family: on the splits we tested, it outperforms the capacity-matched flat Fourier baseline at fewer parameters. The absolute holdout $R^2$ is positive on the smaller split ($N=49$) and negative on the larger ($N=70$), so the method's honest range is bounded by the corpus size. We describe the method, the dataset, and the protocol in full, and we identify four open items that bound the claim.

\section*{References}

\begin{enumerate}
\item Durastanti, C. (2026). \emph{Spectral Bayesian Regression on the Sphere.} arXiv:2601.20528 [math.ST]. \texttt{https://arxiv.org/abs/2601.20528}
\item Anonymous (ICLR 2022 under review). \emph{Generalized Fourier Features for Coordinate-Based Learning on Manifolds.} OpenReview \texttt{g6UqpVislvH}.
\item Sitzmann, V., Martel, J., Bergman, A., Lindell, D., \& Wetzstein, G. (2020). \emph{Implicit Neural Representations with Periodic Activation Functions.} NeurIPS 2020.
\item Tancik, M., et al. (2020). \emph{Fourier Features Let Networks Learn High Frequency Functions in Low Dimensional Domains.} NeurIPS 2020.
\item Mildenhall, B., et al. (2020). \emph{NeRF: Representing Scenes as Neural Radiance Fields for View Synthesis.} ECCV 2020.
\item Cohen, T. S., et al. (2018). \emph{Spherical CNNs.} ICLR 2018.
\item Nickel, M. \& Kiela, D. (2017). \emph{Poincar\'e Embeddings for Learning Hierarchical Representations.} NeurIPS 2017.
\item Ahlfors, L. V. (1979). \emph{Complex Analysis,} 3rd ed. McGraw-Hill.
\item do Carmo, M. P. (1976). \emph{Differential Geometry of Curves and Surfaces.} Prentice-Hall.
\end{enumerate}

\end{document}
